\begin{thebibliography}{}
\bibitem[IEEE]{IEEE}IEEE, {\it The IEEE Standard for Binary Floating-Point Arithmetic}, ANSI/IEEE Std 754-1985. 

\bibitem[ArBo]{ArBo} T. Arbogast, J. L. Bona, Methods of Applied Mathematics, University of Texas at Austin, 2008. Lecture Notes in Applied Mathematics.

\bibitem[Bur] {Bur} A. Burbanks, Renormalization for Siegel disks, Doctoral Thesis, Loughborough University, 1997.

\bibitem[BC]{BC} X. Buff, A. Ch\'eritat, Quadratic Julia sets with positive area, Ann. Math. {\bf 176}:2, (2012), 673-746.

\bibitem[BY]{BY} M. Braverman, M. Yampolsky. Computability of Julia sets. 

\bibitem[dCLM]{CLM} A. de Carvalho, M. Lyubich, M. Martens. Renormalization in the H\'enon family, I: universality but non-rigidity, 
J. Stat. Phys. 121 No. 5/6, (2005), 611-669.

\bibitem[Do] {Do} A. Douady, Disques de Siegel et anneaux de Herman, Seminar Bourbaki 1986/1987, Ast\'erisque \textbf{152-153}, (1987), Soc. Math. France, Montrouge, 1988, 151--172. 

\bibitem[dF1]{dF1} E. de Faria, {\it Proof of universality for critical circle mappings}, Thesis, CUNY, 1992.

\bibitem[dF2]{dF2} E. de Faria, Asymptotic rigidity of scaling ratios for critical circle mappings, {\it Ergodic Theory Dynam. Systems} {\bf 19}(1999), no. 4, 995--1035.

\bibitem[Ga]{Ga}  D. Gaidashev, {\it Cylinder renormalization for Siegel disks and a constructive Measurable Riemann Mapping Theorem},  Nonlinearity {\bf 20}(3) (2007), 713--741. 

\bibitem[GaYaR]{GaRYa} D. Gaidashev, R. Radu, M. Yampolsky, {\it Renormalization and Siegel disks for complex H\'enon maps}, preprint.

\bibitem[GaYa]{GaYa} D. Gaidashev and M. Yampolsky, Cylinder renormalization of Siegel disks, Exp. Math. \textbf{16:2}  (2007), 215--226.


\bibitem[GH] {GH} F. Gehring, K. Hag,  The Ubiquitous Quasidisk, Mathematical Surveys and Monographs \textbf{184} (2012), 171 pp.

\bibitem[Her] {Her} M. Herman, Conjugaison quasi-sym\'etrique des diff\'eomorphismes du cercle et applications auz disques singuliers de Siegel, Manuscripts 1986.

\bibitem[IS]{IS} H. Inou and M. Shishikura. Renormalization for parabolic fixed points and their perturbations.
Preprint. Available from {\sl http://www.math.kyoto-u.ac.jp/$\sim$mitsu/pararenorm/  }

\bibitem[KH]{KH} A. Katok, B. Hasselblatt,  Introduction to the Modern Theory of Dynamical Systems, Cambridge University Press, 1997.

\bibitem[KSW]{KSW} H. Koch, A. Schenkel and P. Wittwer, {\it Computer-assisted proofs in analysis and programming in logic: a case study}, SIAM Rev.   { \bf 38} (1996), 565--604. 

\bibitem[Lan1]{La1} O.E. Lanford, Renormalization group methods for critical circle mappings with general rotation number, {\it VIIIth International Congress on Mathematical Physics (Marseille,1986)}, World Sci. Publishing, Singapore, 532--536, (1987).

\bibitem[Lan2]{La2} O.E. Lanford, Renormalization group methods for critical circle mappings. Nonlinear evolution and chaotic phenomena, {\it NATO adv. Sci. Inst. Ser. B:Phys.} {\bf 176}(1988), Plenum, New York, 25--36.

\bibitem[LM1] {LM1} M. Lyubich, M. Martens. Renormalization of H\'enon maps. Dynamics, Games and Science I, Springer Proceedings in Mathematics 1, M.M. Peixoto, A.A. Pinto,  D.A. Rand (Eds.), (2011), 597-618.

\bibitem[LM2] {LM2}
M. Lyubich, M. Martens. Renormalization in the H\'enon family, II. The 
Heteroclinic Web, Invent. Math. (2001) 186, 115-189.

\bibitem[LM3] {LM3}
M. Lyubich, M. Martens. Probabilistic Universality in Two-Dimensional Dynamics. IMS Stony Brook preprint 11-2.

\bibitem[MP]{MP} R.S. MacKay, I.C. Persival, Universal small-scale structure near the boundary of Siegel disks of arbitrary rotation numer, {\it Physica} {\bf 26D}(1987), 193--202.

\bibitem[MN]{MN} N.S. Manton, M. Nauenberg, Universal scaling behaviour for iterated maps in
the complex plane, {\it Commun. Math. Phys.} {\bf 89}(1983), 555--570.

\bibitem[Mar]{Mar} M. Martens. 
The Periodic Points of Renormalization.
{\it Annals of Math.}
Vol. 147, No. 3 (May, 1998), pp. 543-584

\bibitem[McM1]{McM} C. McMullen, Self-similarity of Siegel disks and Hausdorff dimension of  Julia sets, {\it Acta Math.} {\bf 180}(1998), 247-292.

\bibitem[McM2]{McMBook} C. McMullen,  Renormalization and 3-Manifolds Which Fiber over the Circle (AM-142), Princeton University Press (1996), 253 pages.




%\bibitem[NR]{NR}   W. H. Press, B. P. Flannery, S. A. Teukolsky  and   W. T. Vetterling, {\it Numerical Recipes in Fortran. The Art of Scientific Computing}, Cambridge: Cambridge University Press 1992.

%\bibitem[Shi]{Shish} M. Shishikura, The Hausdorff dimension of the boundary of the Mandelrot set and Julia sets, {\it   Ann. of Math.} {\bf 43}(1942), 607--612.

\bibitem[Sieg]{Sieg} C.L. Siegel, Iteration of analytic functions, {\it Ann. Math.} {\bf 43}(1942), 607-612.

\bibitem[SP]{SP} C. Pugh, M. Shub, Ergodic Attractors, {\it Trans. Amer. Math. Soc.} {\bf 312} 1 (1989), 1-54.

\bibitem[Stir]{Stir} A. Stirnemann,  Existence of the Siegel disc renormalization fixed point,  {\it  Nonlinearity } {\bf 7}(1994),  no. 3, 959--974.


\bibitem[Sw]{Sw} G. \'Swi\c{a}tek, On Critical Circle Homeomorphisms, Bol. Soc. Bras. Mat. \textbf{29:2} (1998), 329-351.

\bibitem[Wi]{Wi} M. Widom, Renormalisation group analysis of quasi-periodicity in analytic maps,
{\it Commun. Math. Phys.} {\bf 92}(1983), 121-136.

\bibitem[Ya1]{Ya1} M. Yampolsky, Hyperbolicity of renormalization of critical circle maps,
{\it  Publ. Math. Inst. Hautes Etudes Sci.} {\bf 96}(2002), 1--41.

\bibitem[Ya2]{Ya2} M. Yampolsky, Renormalization horseshoe for critical circle maps,
{\it  Commun. Math. Physics} {\bf 240}(2003), 75--96.

\bibitem[Ya3]{Ya3} M. Yampolsky,  Siegel disks and renormalization fixed points, Fields Institute Communications, {\bf 53}(2008).

\bibitem[ADA1]{ADA}   S. T. J. Taft and R. A. Duff (eds), {\it Ada 95 Reference Manual: Language and Standard Libraries, International Standard ISO/IEC 8652:1995(E)},  Lec. Notes in Comp. Science {\bf 1246}. 

\bibitem[ADA2]{GNAT}   Ada Core Technologies, 73 Fifth Ave, New York, NY 10003, USA.\\ See also {\tt ftp://cs.nyu.edu/pub/gnat}.

%\bibitem[Prog]{programs} {\tt http://www.math.toronto.edu/gaidash/Programs/siegel-numerics.tar.bz2}.
\end{thebibliography}
